\chapter*{Einleitung zur vierten Lerneinheit}
\addcontentsline{toc}{chapter}{Einleitung zur vierten Lerneinheit}

Viele Aufgaben und Optimierungsprobleme in der Wirtschaft lassen sich durch lineare Zusammenhänge, abhängig von Parametern wie Kapazitäten, Lagerkosten, Fixkosten, Energiebedarf usw. beschreiben. In dieser Lerneinheit werden die mathematischen Grundlagen und Algorithmen zur Lösung solcher Aufgaben vermittelt.

Im ersten Kapitel werden lineare Gleichungssysteme behandelt: Mehrere lineare Gleichungen, die mehrere Unbekannte enthalten, sollen gelöst -- d.\,h. die möglichen Zahlenwerte der Unbekannten sollen bestimmt -- werden. Im Zentrum des Kapitels steht der Gauß-Algorithmus, den die meisten Leser wahrscheinlich bereits im Schulunterricht kennen gelernt (und danach möglicherweise vergessen) haben.

Im mit »Lineare Optimierung« überschriebenen Kapitel wird das Problem behandelt, zu mehreren gegebenen Ungleichungen (ebenfalls mit mehreren Unbekannten) aus den in der Regel unendlich vielen Lösungen (also den möglichen Zahlenwerten für die Unbekannten) solche heraus zu suchen, die hinsichtlich eines gegebenen Kriteriums (»Zielfunktion«) optimal sind. Der hierfür anwendbare Simplex-Algorithmus wird vorgestellt und an Beispielen illustriert.

Wie in den bisherigen Lerneinheiten steht auch hier die Vermittlung der Rechenverfahren (Algorithmen) im Vordergrund, mit denen entsprechende Aufgabenstellungen gelöst werden können. Allerdings kommen auch mathematische Hintergründe zur Sprache, um Zusammenhänge aufzuzeigen und wenigstens eine Ahnung der Möglichkeiten und Grenzen der vorgestellten Verfahren vermitteln zu können.

In Ergänzung der Literaturhinweise in Lerneinheit 1 soll hier noch auf die folgenden zwei Lehrbücher hingewiesen werden, die ebenfalls empfohlen werden können:
\begin{itemize}
  \item Horst Peters: \emph{Wirtschaftsmathematik}. Stuttgart: Kohlhammer Verlag. 4.\ Aufl., 2012
\end{itemize}
